% !TEX program = pdflatex
\documentclass[11pt,a4paper]{article}

\usepackage[margin=1in]{geometry}
\usepackage{amsmath,amssymb}
\usepackage{hyperref}
\usepackage{listings}
\usepackage{xcolor}
\usepackage{enumitem}
\usepackage{booktabs}

\hypersetup{colorlinks=true, linkcolor=blue, urlcolor=blue, citecolor=blue}

\lstdefinelanguage{solidity}{
  keywords={pragma,contract,public,private,external,internal,view,pure,returns,address,uint256,bool,string,mapping,event,modifier,constructor,require,emit,memory,calldata,if,else,for,while,return},
  sensitive=true, comment=[l]{//}, morecomment=[s]{/*}{*/}, morestring=[b]"
}
\lstset{
  basicstyle=\ttfamily\small,
  keywordstyle=\color{blue!70!black}\bfseries,
  commentstyle=\color{green!50!black},
  stringstyle=\color{red!70!black},
  showstringspaces=false,
  breaklines=true,
  frame=single,
  numbers=none
}

\title{\textbf{Trustless Second-Price Auctions on Blockchain}}

\author{Group 7 \;|\; Deming Chen \;\&\; Jieyou Zhao\\ 15-621 Final Project\\
Source code: \url{https://github.com/wodex1nhaoIeng/cmu-15621-final-project}}
\date{\today}

\begin{document}
\maketitle

\begin{abstract}
We present a sealed-bid second-price (Vickrey) auction protocol that runs on
the Ethereum Virtual Machine and removes the need to trust the auctioneer.
In a three-phase commit/encrypt/reveal protocol, bidders publish hash-style
commitments to their bids on-chain, transmit the openings to the seller under
asymmetric encryption, and the seller publishes two non-interactive
zero-knowledge proofs: (i) the winner's commitment opens to the maximum value
of the bid set, and (ii) the publicly announced second-price equals the
opening of some non-winner commitment. Both proofs are constructed from
$\Sigma$-protocols on \texttt{secp256k1} and made non-interactive via
Fiat--Shamir, so they require no trusted setup. We discuss the economic value,
the cryptographic construction, the smart-contract architecture, the threat
model, and the trade-offs of our prototype.
\end{abstract}

% =====================================================================
\section{Economic Problem and Value Proposition}
% =====================================================================

\subsection{What is a second-price auction?}

A \emph{second-price sealed-bid auction} -- also known as the
\emph{Vickrey auction} -- works as follows:

\begin{itemize}[leftmargin=*,nosep]
\item Each bidder submits a sealed (secret) bid.
\item The seller opens all bids simultaneously.
\item The highest bidder wins.
\item The winner pays the \emph{second-highest} price.
\end{itemize}

The mechanism's central appeal is its \textbf{truthful-bidding} property:
each bidder's (weakly) dominant strategy is to bid their true valuation, since
the price they pay is decoupled from the bid they submit. This makes
second-price auctions analytically clean and behaviourally robust.

\paragraph{Worked example.} Suppose Alice bids \$100, Bob \$60, Carol \$80.
Alice wins, but pays only \$80 -- the second-highest bid.

\subsection{The trust problem}

Despite its theoretical beauty, the second-price auction's incentive guarantee
crucially relies on \emph{believing} that the seller will faithfully execute
two steps:
\begin{enumerate}[leftmargin=*,nosep]
\item declare the highest bidder as the winner, and
\item charge \emph{exactly} the second-highest bid.
\end{enumerate}
Neither step is verifiable from a bidder's local view. Two specific attacks
follow immediately, and both are fatal:
\begin{description}[leftmargin=1em,style=nextline]
\item[Attack A: lying about the second price.]
The seller can simply tell the winner that the second-highest bid was higher
than it actually was, charging closer to the winner's full bid.
\item[Attack B: fabricating a phantom bid.]
The seller invents an extra fake bid just below the genuine winner's bid,
extracting nearly the full surplus. In Alice/Bob/Carol's example, the seller
can claim someone bid \$99 and charge Alice \$99 instead of \$80.
\end{description}
There is no way for participants -- bidders or external observers -- to detect
either attack, because the losing bids are never disclosed to anyone but the
seller.

\subsection{Why this matters: market scale}

Second-price logic is not a niche academic curiosity; close
relatives drive a substantial fraction of the modern internet economy. The
\emph{Generalized Second-Price (GSP)} auction, a sponsored-search variant of
Vickrey, was the cash engine of the early-2000s search-advertising boom:
\begin{itemize}[leftmargin=*,nosep]
\item Google's 2005 total revenue was \$6.14\,B; over \textbf{98\%} of it came
from GSP auctions.
\item Yahoo!'s 2005 total revenue was \$5.26\,B, with a large share -- believed
over half -- derived from GSP auctions.
\item As of May 2006 the annual revenues from GSP auctions were on the order
of \textbf{\$10\,B}.
\end{itemize}
Twenty years later the same mechanism family underpins ad markets at Google,
Meta, Amazon, and TikTok, plus large-volume blockchain-native auctions such as
Ethereum Name Service registrations, NFT auctions, and MEV-Boost block
auctions. The trust gap therefore affects markets that move many tens of
billions of dollars annually.

\subsection{Value proposition of a blockchain-native auction}

A blockchain-native second-price auction makes the two-line promise of the
mechanism \emph{cryptographically} enforceable: any participant -- including
third parties -- can verify that the declared winner truly bid the maximum
value and that the price they pay is truly the second-highest. This eliminates
the entire class of post-hoc seller manipulations described above and removes
the need to fall back on reputation, regulation, or law for what is ultimately
a bookkeeping question.

% =====================================================================
\section{How This Is Addressed Today (Non-Blockchain)}
% =====================================================================

The overwhelming majority of second-price-style auctions today run on
\textbf{centralised platforms}.

\paragraph{Pros of centralised platforms.}
\begin{itemize}[leftmargin=*,nosep]
\item Simple to implement and widely understood.
\item Low computational overhead -- it is just a database query.
\end{itemize}

\paragraph{Cons of centralised platforms.}
\begin{itemize}[leftmargin=*,nosep]
\item Bidders \emph{must} trust the platform not to peek, leak, or fabricate bids.
\item There is no way for bidders to verify the second price.
\item The platform is a single point of failure (technical, legal, financial).
\end{itemize}

\paragraph{Representative incumbents.}
\begin{description}[leftmargin=1em,style=nextline]
\item[Google Ads.] Migrated from second-price to first-price in 2019, in part
\emph{because} of trust issues with the second-price mechanism. This is direct
market evidence that the trust gap is real and consequential: a major platform
gave up a theoretically superior mechanism rather than convince advertisers it
was being honest.
\item[eBay.] Uses proxy bidding -- a hidden second-price variant. Users must
trust eBay's internal database; the second price is never independently
verifiable.
\item[Christie's / Sotheby's.] Traditional auction houses where
trustworthiness comes from reputation and regulation rather than from
cryptographic guarantees, and high-profile incidents (chandelier bidding,
ring-bidding scandals) periodically reveal the limits of that approach.
\end{description}
\noindent\textbf{Why these solutions fall short.} Every option above relies on
trust in an intermediary. None offers a cryptographic guarantee that bids were
not manipulated.

% =====================================================================
\section{Our Solution: A Trustless Second-Price Auction on Blockchain}
% =====================================================================

We design a smart-contract-based second-price auction where:
\begin{itemize}[leftmargin=*,nosep]
\item No participant can observe others' bids, even after the auction closes.
\item The seller cannot manipulate or fabricate bids.
\item The correctness of the result is cryptographically verifiable.
\end{itemize}

\subsection{Three-phase protocol}

\begin{description}[leftmargin=1em,style=nextline]
\item[Phase 1 -- Commit (on-chain).]
Each bidder $i$ wishing to bid $p_i$ posts a commitment of the form
$h(p_i \| n_i)$ on-chain, where $n_i$ is a random nonce. In our implementation
$h$ is instantiated by a Pedersen commitment $c_i = p_i\cdot G + n_i\cdot H$
on \texttt{secp256k1}, which is perfectly hiding (so $p_i$ is information-theoretically secret)
and computationally binding under the discrete-logarithm assumption.

\item[Phase 2 -- Reveal to seller (encrypted).]
Bidders encrypt the opening $(p_i, n_i)$ under the seller's public key and
send it to the seller.

\textbf{Design choice: the reveal channel is intentionally left out of the
protocol.} The contract only enforces the on-chain commitment in Phase 1 and
the winner declaration in Phase 3; it deliberately does not prescribe how the
opening reaches the seller. Participants are free to pick whatever channel
fits their threat model -- direct ECIES posted on-chain for public
availability, an off-chain encrypted message, a TLS upload to the seller's
endpoint, etc. This keeps the cryptographic core orthogonal to delivery and
lets the same protocol be reused across very different deployment
environments (public chains, permissioned consortia, mixed on/off-chain
flows). For the demo we instantiate the channel with a Node.js
\texttt{EventEmitter} so that the script can be executed as a single-process
trace, but the seller still cross-checks each reveal against the on-chain
commitment via the \texttt{validReveals} filter, so an inconsistent reveal is
rejected regardless of which channel was used.

\item[Phase 3 -- Settle (zero-knowledge).]
The seller declares the winner on-chain and additionally publishes two
non-interactive zero-knowledge proofs that the winner can verify:
\begin{itemize}[leftmargin=*,nosep]
\item $\pi_{\max}$: the winner's commitment $c_{j^*}$ opens to a value
$p_{j^*}$ such that $p_{j^*} \geq p_i$ for every $i \neq j^*$.
\item $\pi_{\mathrm{mem}}$: the publicly announced second-price $p$ equals
the opening of \emph{some} commitment in
$\{c_i : i \neq j^*\}$, without revealing which one.
\end{itemize}
\end{description}

\subsection{Cryptographic primitives}

We work in the prime-order subgroup of \texttt{secp256k1} with generator $G$
and order $q$. The package \texttt{@noble/curves/secp256k1} supplies fast,
constant-time scalar arithmetic. We derive a second base point
$H = \mathrm{HashToScalar}(\textsf{``internal-pedersen-generator-H''})\cdot G$
whose discrete log relative to $G$ is unknown.

\paragraph{Pedersen commitment.} A bid $p$ with blinding factor
$n \xleftarrow{\$} \mathbb{Z}_q$ commits to
$
c \;=\; p\cdot G \;+\; n\cdot H.
$

\paragraph{Fiat--Shamir.} Every interactive $\Sigma$-protocol below is made
non-interactive by replacing the verifier's challenge with
$
e \;=\; \mathrm{SHA{\text -}256}(\textsf{ctx} \,\|\, \textsf{transcript}) \bmod q,
$
where \textsf{ctx} is a domain-separation string instantiated with the
deployed contract address. This binds proofs to a specific auction and
prevents replay across auctions.

\subsection{$\Sigma$-protocol building blocks}

\begin{enumerate}[leftmargin=*]
\item \textbf{Schnorr discrete-log proof.} Proves knowledge of $w$ with
$Y = w\cdot B$ for any base $B$. Three-move protocol with response
$z = k + e\cdot w$.
\item \textbf{OR-Schnorr.} Given a list of statements $Y_i = w_i\cdot B_i$ the
prover proves knowledge of \emph{some} $w_j$ without revealing the index.
Standard simulation trick: pick random challenges and responses for the
$n-1$ false branches and force their sum to equal the Fiat--Shamir challenge.
\item \textbf{Bit proof.} A commitment $C$ opens to $b\in\{0,1\}$ iff
$
C \;=\; r\cdot H \;\;\text{or}\;\; C - G \;=\; r\cdot H,
$
implemented as a 2-of-2 OR-Schnorr with bases $Y_0 = C$ and $Y_1 = C - G$
relative to $H$.
\item \textbf{Range proof $[0, 2^m)$.} Decompose $v = \sum_{i=0}^{m-1} 2^i b_i$,
commit to each bit, prove each bit is 0/1, and reveal a linear-combination
randomness check. Proof size is $\Theta(m)$.
\end{enumerate}

\subsection{Composite proofs that secure the auction}

\paragraph{Proof of Maximum (\texttt{proveMaximum}).} Given commitments
$c_1,\dots,c_n$ and the prover's index $j^*$, output a proof that
$
\forall\, i\neq j^*:\; p_{j^*} \geq p_i.
$
By the homomorphism of Pedersen, $D_i := c_{j^*} - c_i$ commits to
$p_{j^*} - p_i$ with randomness $n_{j^*} - n_i$. Append a range proof on $D_i$
in $[0, 2^m)$. The bit-length $m$ caps the legal bid range to $[0, 2^m)$,
which (as we show below) is the cornerstone of one of our anti-fraud defences.

\paragraph{Proof of Membership (\texttt{proveMembership}).} Given a public
value $x$ and commitments $c_1,\dots,c_n$, prove that $x = p_{j^\dagger}$ for
some $j^\dagger$ without revealing $j^\dagger$. Reduces to: ``some $c_i - x\cdot G$
opens to $0\cdot G + n_i\cdot H$'', which is an OR-Schnorr with base $H$ over
the $n$ candidate points.

\paragraph{Putting them together.} For an auction with bidders
$\{1,\dots,n\}$ and announced winner index $j^*$, the seller publishes
$j^*$, $\pi_{\max}=\texttt{proveMaximum}(\dots,j^*)$, the second-price $p$,
and $\pi_{\mathrm{mem}}=\texttt{proveMembership}(p,\{c_i\}_{i\neq j^*})$.
The winner verifies (i) $j^* = \texttt{myIndex}$ -- this binds the proof to
the on-chain winner -- (ii) $\pi_{\max}$, and (iii) $\pi_{\mathrm{mem}}$.

\subsection{Smart contract}

The on-chain logic is intentionally minimal:
\begin{lstlisting}[language=solidity]
contract VickreyAuction {
    address public seller;
    uint256 public commitDeadline;
    bool    public finalized;
    address public winner;
    mapping(address => string) public commits; // hex-encoded curve point
    address[] public bidders;

    function commitBid(string calldata c) external; // before deadline
    function finalize(address w) external;          // seller, once
}
\end{lstlisting}
The contract is a public bulletin board for commitments and the single point
of truth for the deadline and the winner identity.

\textbf{Design choice: ZK-proof delivery and verification are intentionally
left to the application layer.} Both $\pi_{\max}$ and $\pi_{\mathrm{mem}}$
are produced by the seller and consumed by whoever needs to be convinced;
in our demo that is the winner, who runs the verifier in JavaScript. The
contract itself only exposes \texttt{commitBid} and \texttt{finalize} and
deliberately stays out of the proof pipeline. This separation of concerns
is valuable for two reasons. First, it lets each deployment pick its own
verification surface: a single suspicious winner can verify locally; an
entire ecosystem can publish the proofs to a watcher service; or, if global
on-chain auditability is desired, the same proofs can be wrapped in a
Groth16 circuit (we ship a generated \texttt{Groth16Verifier.sol}) and
verified by a smart contract. Second, it keeps the on-chain footprint --
and therefore gas cost -- minimal: a single $\pi_{\max}$ is on the order
of kilobytes, and verifying \texttt{secp256k1} arithmetic in the EVM has
no native precompile, so doing it directly on-chain would be prohibitively
expensive. By exposing the proofs as plain data rather than baking the
verifier in, the protocol stays gas-cheap by default while still enabling
full on-chain verification when an application chooses to opt in.

% =====================================================================
\section{Threat Model and Defences}
% =====================================================================

\paragraph{Attack A: seller submits a price just below the highest bid.}
\textbf{Defeated.} Bids are committed in Phase 1 only via hiding hash
commitments. The seller does not learn the highest plaintext bid until Phase 2,
\emph{after} their own commitment window has closed. There is no opportunity
to insert a phantom bid that targets the genuine winner's value.

\paragraph{Attack B: seller submits multiple decoy bids hoping one falls
between the highest and second-highest.}
\textbf{Defeated.} If the seller commits a value higher than the real maximum,
they will be unable to produce $\pi_{\max}$ -- the auction simply fails to
finalise. If the seller commits a value lower than the second-highest, that
decoy is harmless. Either way, the seller cannot move the second price upward.

\paragraph{Attack C: bidder lies on reveal.}
\textbf{Defeated by an explicit cross-check} (\texttt{validReveals} filter
in \texttt{sellerFlow}): the seller recomputes
$\mathrm{commit}(p',n')$ for every reveal and discards mismatches. Without
this filter, a single inconsistent reveal would crash $\pi_{\max}$ generation
and freeze the auction (this was Bug \#4 in our development log).

\paragraph{Attack D: seller swaps the on-chain winner for another address.}
\textbf{Defeated by} the winner-side check $\texttt{proof.maxIndex} =
\texttt{myIndex}$ (Bug \#2 in our log). $\pi_{\max}$ alone only binds
$c_{j^*}$ to the maximum, not $j^*$ to the on-chain winner address.

\paragraph{Attack E: replay across auctions.}
\textbf{Defeated.} Every Fiat--Shamir challenge incorporates the deployed
contract address as \textsf{ctx}, so a proof that verifies against auction $A$
will not verify against auction $B$. Our test suite includes a
``mismatched-ctx'' regression case for both composite proofs.

\paragraph{Attack F: front-running the commit transaction.}
\textbf{Mitigated.} Pedersen commitments are perfectly hiding, so observing
the mempool reveals nothing about the bid value. The bidder need only submit
\texttt{commitBid} before the deadline.

\end{document}
